% Chapter 9 of Optimizers in the Wild: warmup.

\casechapter{Warmup}{Why the step size must start small}{Linear warmup}{ch:warmup}
\begin{refsection}[chapters/09-warmup.bib]

\begin{chaptersummary}
Large training runs raise the learning rate from nearly zero to its peak over the first thousands of steps. The smallest idealization that makes this necessary bounds the curvature by an affine function of the loss gap, $S\le K_0+K_1(\Loss-\Loss^\star)$, so that the step that is safe at a point grows as the loss falls. We prove a descent lemma for this class and show that, under the Polyak--Łojasiewicz inequality, a warmup indexed by the loss reaches accuracy $\eps$ in $(K_0\log(\Delta_0/\eps)+K_1\Delta_0)/\mu$ steps while a constant step pays the initial curvature again in every factor of $\eps$; the optimal profile is a rising hyperbola that saturates, given by the Wright omega function, and a model that saturates both hypotheses shows it optimal among stable schedules and a linear ramp too slow by the factor $u_0/(4e)$. We also prove that Adam's first step is a sign step and that its second moment averages only the few gradients seen so far, and we derive that a warmup at the edge of stability drives the sharpness down along $2/\eta_t$, as a small network confirms. The theory does not derive its curvature bound for networks, where it fails globally, and it misses the catapults through which shorter, unstable warmups succeed.
\end{chaptersummary}

\section{In the wild: warming up the learning rate}\label{sec:warmup:wild}

The learning rate of a modern training run is a curve, and the curve starts low. An early instance in deep learning is in the paper that introduced residual networks: to train a network with $110$ layers on CIFAR-10, \textcite{HeEtAl2016} used the learning rate $0.01$ until the training error fell below $80\%$, about $400$ iterations, and only then the rate $0.1$ of their other networks, which they found too large to start converging.\sidenote{In this chapter \emph{learning rate}, the word of the practice sections, and \emph{step size}, the word of the theory, mean the same thing.} \Textcite{GoyalEtAl2017} made the ramp gradual. To train a ResNet-50 on ImageNet with minibatches of $8192$ images, they multiplied the learning rate by the factor by which the minibatch had grown, the linear scaling rule of \cref{ch:noise}, and reached that rate by a linear increase over the first five epochs, because the rule fails early in training, when the network changes rapidly. The original Transformer \parencite{VaswaniEtAl2017} increased its learning rate linearly for $4000$ steps and then decayed it as the inverse square root of the step. Large language models inherited the ramp.

\begin{practice}[GPT-3, \textcite{BrownEtAl2020}; LLaMA, \textcite{TouvronEtAl2023}; Adam's rule of thumb, \textcite{MaYarats2021}]
Raise the learning rate linearly from (nearly) zero to its peak, over the first $375$ million tokens for GPT-3 and the first $2000$ steps for LLaMA, then decay it along a cosine to $10\%$ of the peak (\cref{ch:schedules}); both runs use AdamW with $\beta_2=0.95$. For Adam a rule of thumb sets the length of a linear warmup to $2/(1-\beta_2)$ steps: $2000$ steps for $\beta_2=0.999$, and $40$ for $\beta_2=0.95$. Otherwise the length is inherited from an earlier recipe or tuned by a sweep; we know of no derivation of these numbers.
\end{practice}

What has been observed about warmup falls into three groups.

\begin{observation}[Warmup and the placement of layer normalization; \textcite{XiongEtAl2020}]\label[observation]{obs:warmup:postln}
In the original Transformer, which normalizes after each residual addition (Post-LN), the gradients of the parameters near the output are large at initialization, so that a large initial learning rate makes training unstable, and the warmup stage avoids this. With the normalization inside the residual branches (Pre-LN) the gradients are well behaved at initialization, and the warmup can be removed: Pre-LN Transformers trained without it reach results comparable to the baselines with less training time and less tuning.
\end{observation}

The statement about the gradients is a theorem of \textcite{XiongEtAl2020}, proved with mean-field approximations at initialization; the rest is experimental, on translation and language-model pretraining. An initialization can remove the need as well: \textcite{HuangEtAl2020} train Transformers without warmup and without layer normalization by rescaling their initial weights.

\begin{observation}[Warmup lowers the sharpness; \textcite{GilmerEtAl2022}, \textcite{KalraBarkeshli2024}]\label[observation]{obs:warmup:sharpness}
Early instabilities coincide with a sharpness well above $2/\eta$, and warmup stabilizes training by letting the sharpness fall before the learning rate is high \parencite{GilmerEtAl2022}. When a network starts in a region sharper than the target learning rate tolerates, its sharpness falls during warmup and stays near $2/\eta_t$, and the main benefit of warmup is that larger target learning rates become trainable \parencite{KalraBarkeshli2024}.
\end{observation}

The \term{sharpness}\index{sharpness} is the largest eigenvalue of the Hessian, $S(\theta)=\lmax(\nabla^2\Loss(\theta))$, and $2/\eta$ is the stability threshold of gradient descent on a quadratic (\cref{ch:eos}). \Cref{fig:warmup:mlp} shows the second pattern in a network small enough to train in seconds: the sharpness falls from $132$ to $20$ while the learning rate rises to $0.1$, and without warmup the same learning rate diverges in six steps.

\begin{observation}[Explanations specific to Adam]\label[observation]{obs:warmup:adam}
\textcite{LiuEtAl2020RAdam} attributed Adam's need for warmup to the variance of its adaptive learning rate in the first steps, when the second-moment estimate has averaged few gradients, and proposed RAdam, which rectifies the step until that variance is tractable. \textcite{MaYarats2021} found that an untuned linear warmup over $2/(1-\beta_2)$ steps works as well as RAdam. \textcite{KossonMessmerJaggi2024} attributed the need for warmup in GPT training to large angular updates of the weights in the first steps, and reduced it by controlling them. \textcite{KalraBarkeshli2024} initialized Adam's second moment with the first squared gradient (GI-Adam) and reported benefits similar to those of warmup.
\end{observation}

The three groups are not rival theories. The first two are about curvature, the third about Adam's statistics and the size of its early updates, and a transformer trained with Adam has all of them. What is robust, across architectures and optimizers, is that some warmup prevents early divergence at the learning rates that practice wants; in small transformers a longer warmup also makes the final loss less sensitive to the learning rate \parencite{WortsmanEtAl2024}. What is folklore is the length. In \cref{obs:warmup:sharpness} the curvature falls because the schedule rises. The idealization of this chapter starts from the other end, a curvature that falls because the loss does, and asks two questions: why must the step size start small, and how should it grow?

\section{The idealization: curvature bounded by the loss}\label{sec:warmup:ideal}

Gradient descent on a function whose curvature is at most $L$ everywhere decreases the loss at every step if $\eta<2/L$ (\cref{lem:toolkit:descent}). A network has no useful global $L$: the sharpness of the network of \cref{fig:warmup:mlp} is $132$ at initialization and $20$ after a few hundred steps, and the sharpness of a network with a rescaling symmetry is not even bounded on its set of minimizers (\cref{rem:warmup:symmetry}). The idealization replaces the global bound by one that improves as the loss falls.

\begin{assumption}[Loss-dependent curvature]\label[assumption]{ass:warmup:curvature}
The loss $\Loss\colon\R^d\to\R$ is twice continuously differentiable and bounded below, with $\Loss^\star=\inf\Loss$ and gap $\Delta(\theta)=\Loss(\theta)-\Loss^\star$. For the initial point $\theta_0$ let $\Theta_0=\{\theta\colon\Loss(\theta)\le\Loss(\theta_0)\}$. There are constants $K_0,K_1\ge0$ such that
\begin{equation}\label{eq:warmup:curvature}
  S(\theta)\le K_0+K_1\Delta(\theta)\qquad\text{for every }\theta\in\Theta_0 .
\end{equation}
\end{assumption}

The \term{loss-dependent curvature}\index{loss-dependent curvature} condition gives the curvature $K_0$ that may remain at the bottom and the excess $K_1\Delta$ that a gap $\Delta$ allows. It restricts the Hessian from above only; nothing in this chapter uses a lower bound. \Textcite{AlimisisIslamovLucchi2026} introduced the condition with the operator norm $\norm{\nabla^2\Loss}_{\op}$ in place of $S$, which implies \eqref{eq:warmup:curvature}, and \textcite{RiabininEtAl2026} state it in a general norm with an extra term $K_\rho\Delta^\rho$ (\cref{sec:warmup:normalized}).\sidenote{Its relative for the gradient, $\norm{\nabla^2\Loss}\le L_0+L_1\norm{\nabla\Loss}$, is the $(L_0,L_1)$-smoothness of \textcite{ZhangHeSraJadbabaie2020} and of \cref{ch:spikes}. \Textcite{AlimisisIslamovLucchi2026} prefer the loss to the gradient because the gradient norm can rise early in training, whereas the loss falls along safe steps.} A quadratic satisfies it with $K_1=0$. The function
\begin{equation}\label{eq:warmup:cosh}
  \Loss_{\mathrm c}(z)=\frac{K_0}{K_1}\Bigl(\cosh\bigl(\sqrt{K_1}\,z\bigr)-1\Bigr),\qquad z\in\R,
\end{equation}
satisfies it with equality, $\Loss_{\mathrm c}''=K_0\cosh(\sqrt{K_1}z)=K_0+K_1\Loss_{\mathrm c}$: its curvature grows as fast as its loss allows, and $2\cosh z-2$ is the case $K_0=2$, $K_1=1$. The second hypothesis says that the gradient is large when the loss is.

\begin{assumption}[Polyak--Łojasiewicz inequality]\label[assumption]{ass:warmup:pl}
There is $\mu>0$ such that $\norm{\nabla\Loss(\theta)}_2^2\ge2\mu\Delta(\theta)$ for every $\theta\in\Theta_0$.
\end{assumption}

Under the \term{Polyak--Łojasiewicz inequality}\index{Polyak--Łojasiewicz inequality} every stationary point in $\Theta_0$ is a global minimizer, and with a global curvature bound $L$ gradient descent converges at the linear rate $1-\mu/L$ (\cref{lem:toolkit:pl}).\sidenote{The inequality is due to Polyak and to Łojasiewicz, both in 1963; see \textcite{Polyak1987} and, for its uses in machine learning, \textcite{KarimiNutiniSchmidt2016}.} Both hypotheses are statements about the loss, and that is the point of combining them: the first says how large a step the current loss allows, the second how much a step of that size achieves.

The extremal member of the class has two kinds of coordinates.

\begin{definition}[The warmup model]\label[definition]{def:warmup:model}
Fix $\mu>0$, $K_0\ge2\mu$ and $K_1>0$. The \term{warmup model}\index{warmup model} is the loss
\[
  \Loss(x,z)=q(x)+\tfrac12S(x)\norm{z}_2^2,\qquad q(x)=\tfrac\mu2\norm{x}_2^2,\qquad S(x)=K_0+K_1q(x),
\]
on $\theta=(x,z)\in\R^{d_x}\times\R^{d_z}$, trained by gradient descent from a point $(x_0,z_0)$ with $z_0$ small.
\end{definition}

The coordinate $x$ is slow: it carries the loss, with curvature $\mu$. The coordinate $z$ is stiff: its curvature $S(x)$ is set by the loss of the slow coordinate. The model caricatures the start of training, where the loss is high because the network has not yet fitted the data, and the curvature in the stiffest direction is high for the same reason.

\begin{proposition}[The warmup model saturates both hypotheses]\label[proposition]{prop:warmup:model}
In the warmup model $\Loss^\star=0$, and the following hold.
\begin{enumerate}
\item \Cref{ass:warmup:curvature,ass:warmup:pl} hold on all of $\R^{d_x+d_z}$ with the constants $K_0$, $K_1$ and $\mu$, and even $\norm{\nabla^2\Loss(\theta)}_{\op}\le K_0+K_1\Delta(\theta)$.
\item On the subspace $\{z=0\}$ both hold with equality: $S=K_0+K_1\Delta$ and $\norm{\nabla\Loss}_2^2=2\mu\Delta$.
\item Gradient descent with steps $\eta_t$ satisfies
\begin{equation}\label{eq:warmup:modelmap}
  x_{t+1}=\Bigl(1-\eta_t\mu\bigl(1+\tfrac12K_1\norm{z_t}_2^2\bigr)\Bigr)x_t,\qquad z_{t+1}=\bigl(1-\eta_tS(x_t)\bigr)z_t .
\end{equation}
So the subspace $\{z=0\}$ is invariant, on it $\Delta_{t+1}=(1-\mu\eta_t)^2\Delta_t$, and a perturbation of the stiff coordinate is amplified at step $t$ exactly when $\eta_tS(x_t)>2$.
\end{enumerate}
\end{proposition}

\begin{proof}
Write $\varepsilon=\frac12K_1\norm{z}_2^2$. Then $\nabla_x\Loss=\mu(1+\varepsilon)x$, $\nabla_z\Loss=S(x)z$, $\Loss=q(1+\varepsilon)+\frac12K_0\norm{z}_2^2$ and $K_0+K_1\Loss=S(1+\varepsilon)$. Part (iii) follows from the gradient, and part (ii) from setting $z=0$, where the Hessian is block diagonal with blocks $\mu I$ and $S(x)I$, and $S(x)\ge K_0\ge\mu$.

For (i), the Hessian has the diagonal blocks $\mu(1+\varepsilon)I$ and $SI$ and the off-diagonal block $K_1\mu\,xz^\top$. On vectors orthogonal to $x$ in the first block or to $z$ in the second it acts as $\mu(1+\varepsilon)$ or as $S$, and both lie in $[0,S(1+\varepsilon)]$. If $x$ and $z$ are nonzero, it remains to bound the $2\times2$ matrix $A$ that it induces on the span of $(x/\norm{x}_2,0)$ and $(0,z/\norm{z}_2)$, with diagonal entries $\mu(1+\varepsilon)$ and $S$ and off-diagonal entry $b=K_1\mu\norm{x}_2\norm{z}_2$, for which $b^2=4\mu K_1q\varepsilon$. The matrix $S(1+\varepsilon)I-A$ has the nonnegative diagonal entries $(S-\mu)(1+\varepsilon)$ and $S\varepsilon$, and its determinant is at least $(S-\mu)S\varepsilon-b^2=\varepsilon\bigl((S-\mu)S-4\mu K_1q\bigr)$. The bracket is nonnegative: $K_0\ge2\mu$ gives $S-\mu\ge K_0/2+K_1q$, and the inequality of arithmetic and geometric means gives $(K_0/2+K_1q)(K_0+K_1q)\ge2\sqrt2\,K_0K_1q\ge4\mu K_1q$. Hence $\lmax(A)\le S(1+\varepsilon)=K_0+K_1\Loss$. Also $\lmin(A)\ge-b$, and $b=2\sqrt{K_1q\cdot\mu\varepsilon}\le K_1q+\mu\varepsilon\le S(1+\varepsilon)$, which gives the bound on the operator norm. Finally, because $S^2\ge K_0^2\ge\mu K_0$,
\[
  \norm{\nabla\Loss}_2^2=2\mu q(1+\varepsilon)^2+S^2\norm{z}_2^2\ge2\mu q(1+\varepsilon)+\mu K_0\norm{z}_2^2=2\mu\Loss .\qedhere
\]
\end{proof}

So on the line $z=0$ the warmup model has the most curvature that \cref{ass:warmup:curvature} allows and the least gradient that \cref{ass:warmup:pl} allows. Gradient descent on it is a scalar recursion as long as the stiff coordinate stays quiet, and the instability of that coordinate is the only way in which a step can fail. The hypothesis $K_0\ge2\mu$ is a convenience (Exercise~\ref{ex:warmup:threshold}).

\begin{dictionary}
the gap $\Delta_t$ & the training loss minus its floor, which is unknown in practice\\
$S\le K_0+K_1\Delta$ & the sharpness falls as the loss falls early in training (measured, not derived)\\
$K_0$, and the final step $1/K_0$ & the curvature late in training, which limits the peak learning rate\\
the excess $u_0=K_1\Delta_0/K_0$ & how much sharper the initialization is than the trained network\\
the PL constant $\mu$, and $\kappa=K_0/\mu$ & the slowest direction of descent, and a condition number\\
the loss-indexed step $1/(K_0+K_1\Delta_t)$ & a warmup that follows the loss rather than the step count\\
slow coordinate $x$, stiff coordinate $z$ & the bulk of the parameters, and the top eigenvector of the Hessian\\
$\eta_tS_t>2$ & the edge of stability (\cref{ch:eos}): a catapult or a divergence\\
Adam's $n_t$ & the number of gradients that the second-moment estimate has effectively averaged
\end{dictionary}

\paragraph{What was thrown away.} Noise: the gradients are exact, so the gradient noise behind the warmup of \textcite{GoyalEtAl2017} is absent (\cref{ch:noise}). The response of the landscape: \cref{ass:warmup:curvature} is a fixed inequality, whereas in a network the sharpness responds to the step size (\cref{ch:eos}). The geometry: in a network the Polyak--Łojasiewicz inequality can hold at best near the trajectory, and \cref{ass:warmup:curvature} only on a bounded region (\cref{rem:warmup:symmetry}). The optimizer: gradient descent without momentum or preconditioning, until \cref{sec:warmup:more}. Each simplification removes a mechanism that practice may rely on, so that what remains is the one that the hypothesis isolates: the curvature at the start exceeds what the target step tolerates, and it falls as training proceeds. If warmup were not needed in this idealization, it would not be needed for this reason in a network either.

\paragraph{What one had to notice.} The safe step is a property of the current point, and it can be indexed by the loss, which, unlike the gradient norm, does not increase along safe steps. A step that is safe at the current loss is therefore safe at every later one, a bound in terms of the loss turns into a bound in terms of time, and an increasing schedule falls out of it.

\section{The theory: a safe step that grows as the loss falls}\label{sec:warmup:theory}

The classical descent lemma integrates a curvature bound along the segment of a step. Under \cref{ass:warmup:curvature} the bound is available only where the loss is below its starting value, and the segment is not known in advance to stay there. A continuity argument closes the gap.

\begin{lemma}[A local descent lemma]\label[lemma]{lem:warmup:descent}
Let $\Loss$ be twice continuously differentiable, let $\theta,u\in\R^d$ satisfy $\ip{\nabla\Loss(\theta),u}<0$, and let $M>0$ be such that $u^\top\nabla^2\Loss(\theta+su)\,u\le M$ for every $s\ge0$ with $\Loss(\theta+su)\le\Loss(\theta)$. Then for every $s\in\bigl[0,2\abs{\ip{\nabla\Loss(\theta),u}}/M\bigr]$ the segment from $\theta$ to $\theta+su$ lies in $\{\Loss\le\Loss(\theta)\}$, and
\[
  \Loss(\theta+su)\le\Loss(\theta)+s\ip{\nabla\Loss(\theta),u}+\tfrac12Ms^2 .
\]
\end{lemma}

\begin{proof}
Let $\varphi(r)=\Loss(\theta+ru)$, $\gamma=-\varphi'(0)=\abs{\ip{\nabla\Loss(\theta),u}}$ and $\bar s=2\gamma/M$, and fix $s\in(0,\bar s]$; the case $s=0$ is trivial. The plan is to show that $\varphi$ cannot climb back to $\varphi(0)$ before $s$. Let $A$ be the set of $r\in[0,s]$ such that $\varphi\le\varphi(0)$ on $[0,r]$. It contains $0$, it is an interval, and it is closed because $\varphi$ is continuous, so $A=[0,r^*]$. On $A$ the hypothesis gives $\varphi''\le M$, and Taylor's formula with integral remainder gives, for $r\in A$,
\begin{equation}\label{eq:warmup:taylor}
  \varphi(r)\le\varphi(0)-\gamma r+\tfrac12Mr^2=\varphi(0)-\gamma r\bigl(1-r/\bar s\bigr).
\end{equation}
Suppose that $r^*<s$. If $r^*=0$, then $\varphi'(0)<0$ gives $\varphi<\varphi(0)$ on some interval $(0,\delta)$. If $r^*>0$, then $0<r^*<\bar s$ and \eqref{eq:warmup:taylor} give $\varphi(r^*)<\varphi(0)$, and by continuity $\varphi<\varphi(0)$ on some interval $[r^*,r^*+\delta)$. Either way $A$ contains points beyond $r^*$, which is absurd. So $r^*=s$, the segment lies in the sublevel set, and \eqref{eq:warmup:taylor} at $r=s$ is the claim.
\end{proof}

For gradient descent at a point $\theta\in\Theta_0$ with gap $\Delta$ and gradient $g\ne0$, take $u=-g$. The points of the ray with $\Loss\le\Loss(\theta)$ lie in $\Theta_0$ and have gap at most $\Delta$, so \cref{ass:warmup:curvature} gives the hypothesis of the lemma with $M=(K_0+K_1\Delta)\norm{g}_2^2$. Hence, for $0\le\eta\le2/(K_0+K_1\Delta)$,
\begin{equation}\label{eq:warmup:gd}
  \Loss(\theta-\eta g)\le\Loss(\theta)-\eta\Bigl(1-\frac{\eta(K_0+K_1\Delta)}2\Bigr)\norm{g}_2^2,
\end{equation}
and in particular $\Loss(\theta-\eta g)\le\Loss(\theta)-\frac\eta2\norm{g}_2^2$ if $\eta\le1/(K_0+K_1\Delta)$. This is the descent lemma with the global constant replaced by the curvature that the current loss allows. The step does not increase the loss, so the step that \eqref{eq:warmup:gd} certifies at the next point is at least as large: the \term{local descent lemma}\index{local descent lemma} already contains a warmup.

\begin{theorem}[Warmup under loss-dependent curvature]\label{thm:warmup:schedule}
Let \cref{ass:warmup:curvature,ass:warmup:pl} hold with $K_0>0$, let $\Delta_0=\Delta(\theta_0)>0$, and define $r(D)=\mu D/(K_0+K_1D)$ and the scalar sequence
\[
  D_0=\Delta_0,\qquad D_{t+1}=D_t-r(D_t).
\]
Then $\mu\le K_0$, and the gaps $\Delta_t=\Delta(\theta_t)$ of gradient descent satisfy the following.
\begin{enumerate}
\item \emph{Loss-indexed warmup.} With $\eta_t=1/(K_0+K_1\Delta_t)$, $\Delta_{t+1}\le\Delta_t-r(\Delta_t)$ for every $t$.
\item \emph{Scheduled warmup.} With $\eta_t=1/(K_0+K_1D_t)$, $\Delta_t\le D_t$ for every $t$. This schedule is nondecreasing, from $1/(K_0+K_1\Delta_0)$ to the limit $1/K_0$.
\item In both cases, for every $t\ge0$,
\begin{equation}\label{eq:warmup:time}
  K_0\log\frac{\Delta_0}{\Delta_t}+K_1(\Delta_0-\Delta_t)\ge\mu t,
\end{equation}
so that $\Delta_T\le\eps$ for every $\eps\in(0,\Delta_0]$ and every $T\ge\bigl(K_0\log(\Delta_0/\eps)+K_1(\Delta_0-\eps)\bigr)/\mu$.
\item \emph{Constant step.} With $\eta_t=\eta\le1/(K_0+K_1\Delta_0)$, $\Delta_t\le(1-\mu\eta)^t\Delta_0$; at the largest such $\eta$ this guarantees $\Delta_T\le\eps$ for $T\ge(K_0+K_1\Delta_0)\log(\Delta_0/\eps)/\mu$.
\end{enumerate}
\end{theorem}

\begin{proof}
The plan is to apply \eqref{eq:warmup:gd} at every step, after checking that the step is admissible there, and to turn the one-step decrease into the time bound \eqref{eq:warmup:time} by comparison with an integral. A step from a point with $\Delta_t=0$ does not move, because such a point is a global minimizer; and at a point of $\Theta_0$ with $\Delta_t>0$ the gradient is nonzero by \cref{ass:warmup:pl}. We use the convention $\log(\Delta_0/0)=+\infty$.

\emph{Part (i).} Suppose that $\theta_t\in\Theta_0$ and $\Delta_t>0$. The step $\eta_t=1/(K_0+K_1\Delta_t)$ is admissible in \eqref{eq:warmup:gd}, and with \cref{ass:warmup:pl},
\[
  \Delta_{t+1}\le\Delta_t-\frac{\norm{\nabla\Loss(\theta_t)}_2^2}{2(K_0+K_1\Delta_t)}\le\Delta_t-\frac{\mu\Delta_t}{K_0+K_1\Delta_t}=\Delta_t-r(\Delta_t).
\]
In particular $\theta_{t+1}\in\Theta_0$, and (i) follows by induction.

\emph{The inequality $\mu\le K_0$.} The same computation at an arbitrary point of $\Theta_0$ with positive gap $\Delta$, together with $\Loss^\star\le\Loss(\theta-\eta g)$, gives $\mu\le K_0+K_1\Delta$ there. The iterates of (i) have $\Delta_{t+1}\le(1-\mu/(K_0+K_1\Delta_0))\Delta_t$, so their gaps tend to $0$. If a gap vanishes, at a first step $T$, then by \cref{lem:warmup:descent} the segment from $\theta_{T-1}$ to $\theta_T$ lies in $\Theta_0$, and $\Loss$ takes on it every value between $\Loss^\star$ and $\Loss(\theta_{T-1})$. Either way $\Theta_0$ contains points of arbitrarily small positive gap, and $\mu\le K_0$.

\emph{Part (ii).} Because $\mu\le K_0$, the map $D\mapsto D-r(D)=D\bigl(1-\mu/(K_0+K_1D)\bigr)$ sends $[0,\infty)$ into itself, so the $D_t$ are nonnegative and nonincreasing; they tend to $0$ because $D_{t+1}\le(1-\mu/(K_0+K_1\Delta_0))D_t$. Hence $\eta_t$ is nondecreasing with limit $1/K_0$. For the comparison, suppose that $\Delta_t\le D_t$, as holds at $t=0$. If $\Delta_t=0$ there is nothing to prove. Otherwise $\eta_t=1/(K_0+K_1D_t)\le1/(K_0+K_1\Delta_t)$ is admissible in \eqref{eq:warmup:gd}, and
\[
  \Delta_{t+1}\le\Delta_t-\frac{\eta_t}2\norm{\nabla\Loss(\theta_t)}_2^2\le(1-\mu\eta_t)\Delta_t\le(1-\mu\eta_t)D_t=D_{t+1},
\]
where the last inequality uses $1-\mu\eta_t\ge1-\mu/K_0\ge0$.

\emph{Part (iii).} Let $\Phi(D)=K_0\log D+K_1D$ for $D>0$ and $\Phi(0)=-\infty$. It is increasing, and its derivative $\Phi'(D)=(K_0+K_1D)/D$ is decreasing. Let $(E_t)$ be a nonnegative sequence with $E_0=\Delta_0$ and $E_{t+1}\le E_t-r(E_t)$, such as the gaps of (i) or the numbers $D_t$ of (ii). If $E_{t+1}>0$, then, because $\Phi'$ is smallest at the right end of $[E_{t+1},E_t]$,
\[
  \Phi(E_t)-\Phi(E_{t+1})=\int_{E_{t+1}}^{E_t}\Phi'(s)\,\dd s\ge(E_t-E_{t+1})\,\Phi'(E_t)\ge r(E_t)\,\frac{K_0+K_1E_t}{E_t}=\mu .
\]
Summing, $\Phi(E_t)\le\Phi(\Delta_0)-\mu t$, which also holds when some $E_s$ vanishes. In case (ii), $\Delta_t\le D_t$ and the monotonicity of $\Phi$ give $\Phi(\Delta_t)\le\Phi(D_t)$. So in both cases $\Phi(\Delta_t)\le\Phi(\Delta_0)-\mu t$, which is \eqref{eq:warmup:time}. If $T\ge(\Phi(\Delta_0)-\Phi(\eps))/\mu$, then $\Phi(\Delta_T)\le\Phi(\eps)$, so $\Delta_T\le\eps$.

\emph{Part (iv).} If $\Delta_t\le\Delta_0$, then $\eta\le1/(K_0+K_1\Delta_t)$ is admissible in \eqref{eq:warmup:gd}, which gives $\Delta_{t+1}\le(1-\mu\eta)\Delta_t$. By induction $\Delta_t\le(1-\mu\eta)^t\Delta_0\le e^{-\mu\eta t}\Delta_0$, which is at most $\eps$ once $t\ge\log(\Delta_0/\eps)/(\mu\eta)$.
\end{proof}

Parts (i) to (iii) are the warmup: a step size that starts at the value $1/(K_0+K_1\Delta_0)$ that the initial curvature allows and rises to $1/K_0$ as the loss falls. The schedule of (ii) is a fixed function of the step count, computed in advance from four numbers. The rule of (i) needs the gap, which in practice is unknown, because $\Loss^\star$ is. The bound \eqref{eq:warmup:time} is additive: the initial excess curvature costs $K_1\Delta_0/\mu$ steps once, and the final curvature $K_0/\mu$ steps per factor $e$ of accuracy. The constant step of (iv) pays $(K_0+K_1\Delta_0)/\mu$ steps per factor $e$. The shape of the schedule can be computed exactly.

\begin{corollary}[The warmup profile]\label[corollary]{cor:warmup:profile}
In the setting of \cref{thm:warmup:schedule} let $K_1>0$ and $u_0=K_1\Delta_0/K_0$, and for real $s$ let $\omega(s)$ be the positive solution $w$ of $w+\log w=s$. The function
\[
  D(t)=\frac{K_0}{K_1}\,\omega\Bigl(u_0+\log u_0-\frac{\mu t}{K_0}\Bigr),\qquad t\ge0,
\]
solves $\dot D=-r(D)$ with $D(0)=\Delta_0$, and the numbers $D_t$ of the scheduled warmup satisfy, if $\mu<K_0$,
\begin{equation}\label{eq:warmup:sandwich}
  D\Bigl(\frac{t}{1-\mu/K_0}\Bigr)\le D_t\le D(t)\qquad\text{for every }t\ge0
\end{equation}
(the right inequality also when $\mu=K_0$). So the scheduled step $1/(K_0+K_1D_t)$ lies between the \term{warmup profile}
\[
  \eta(t)=\frac{1}{K_0\bigl(1+\omega(u_0+\log u_0-\mu t/K_0)\bigr)}
\]
and the same profile at time $t/(1-\mu/K_0)$. If $u_0\ge1$, the profile reaches half its final value, $1/(2K_0)$, at the time $T_{1/2}=(K_0/\mu)(u_0-1+\log u_0)$.
\end{corollary}

\begin{proof}
With $u=K_1D/K_0$, the function $\Phi$ of the proof of \cref{thm:warmup:schedule} is $\Phi(D)=K_0(u+\log u)+K_0\log(K_0/K_1)$, so $D(t)$ is defined by $\Phi(D(t))=\Phi(\Delta_0)-\mu t$. Differentiating, $\Phi'(D)\dot D=-\mu$, that is, $\dot D=-\mu D/(K_0+K_1D)=-r(D)$. The right inequality of \eqref{eq:warmup:sandwich} is $\Phi(D_t)\le\Phi(\Delta_0)-\mu t=\Phi(D(t))$, proved in part (iii) of the theorem, with the monotonicity of $\Phi$. For the left one, $\Phi'$ is largest at the left end of $[D_{t+1},D_t]$, and $D_{t+1}/D_t=1-\mu/(K_0+K_1D_t)\ge1-\mu/K_0$, so
\[
  \Phi(D_t)-\Phi(D_{t+1})\le(D_t-D_{t+1})\frac{K_0+K_1D_{t+1}}{D_{t+1}}\le r(D_t)\,\frac{K_0+K_1D_t}{D_{t+1}}=\mu\frac{D_t}{D_{t+1}}\le\frac{\mu}{1-\mu/K_0} .
\]
Summing, $\Phi(D_t)\ge\Phi(\Delta_0)-\mu t/(1-\mu/K_0)=\Phi\bigl(D(t/(1-\mu/K_0))\bigr)$. The time $T_{1/2}$ is the solution of $\omega(u_0+\log u_0-\mu t/K_0)=1$, that is, of $1=u_0+\log u_0-\mu t/K_0$.
\end{proof}

The function $\omega$ is the Wright omega function, $\omega(s)=W(e^s)$ with $W$ the principal branch of the Lambert function.\sidenote{The Lambert function inverts $w\mapsto we^w$, so $w+\log w=s$ means $we^w=e^s$. It is available as \texttt{scipy.special.lambertw}, and $\omega$ as \texttt{scipy.special.wrightomega}.} The profile has two phases. While $u=K_1D/K_0$ is large, $u+\log u=u_0+\log u_0-\mu t/K_0$ gives $u\approx u_0-\mu t/K_0$: the gap falls by about $\mu/K_1$ per step, and the step size rises along the hyperbola $\eta(t)\approx1/(K_0+K_1\Delta_0-\mu t)$, slowly at first and then faster. Once $u$ is small, $u\approx u_0e^{u_0}e^{-\mu t/K_0}$: the gap decays at the rate $\mu/K_0$, and the step size approaches $1/K_0$ exponentially fast. The profile is convex while $\eta<2/(3K_0)$ and concave afterwards (Exercise~\ref{ex:warmup:profile}). The \term{warmup}\index{warmup!length} is the first phase, and its length $T_{1/2}$ is the initial excess curvature $K_1\Delta_0$ divided by $\mu$, up to a logarithm. The sandwich \eqref{eq:warmup:sandwich} says that the continuum describes the discrete schedule up to a relative error $1/\kappa$ in time, where $\kappa=K_0/\mu$.

How much does the warmup buy? A lower bound for constant steps needs a hard function, and the cosh function \eqref{eq:warmup:cosh} supplies one: a constant step large enough for the late phase overshoots it at the start, while a quadratic punishes every step small enough for the start.

\begin{proposition}[A constant step pays the initial curvature at every scale]\label[proposition]{prop:warmup:constant}
Let $0<\mu\le K_0$, $K_1>0$, $\Delta_0>0$ and $0<\eps<\Delta_0/2$, and let
\[
  \Loss(x,z)=\tfrac\mu2x^2+\Loss_{\mathrm c}(z),\qquad(x,z)\in\R^2,
\]
with $\Loss_{\mathrm c}$ as in \eqref{eq:warmup:cosh}, started at a point $(x_0,z_0)$ with $\frac\mu2x_0^2=\Loss_{\mathrm c}(z_0)=\Delta_0/2$ and $z_0>0$. Then $\Loss^\star=0$, \cref{ass:warmup:curvature,ass:warmup:pl} hold on $\R^2$ with the constants $K_0$, $K_1$ and $\mu$, and with $a=\sqrt{K_1}z_0$ and $\eta^*=2a/(K_0\sinh a)$ gradient descent with a constant step $\eta$
\begin{enumerate}
\item never reaches $\Delta_t\le\eps$ if $\eta\ge\eta^*$, and
\item needs at least $(1-\eta\mu)\log(\Delta_0/(2\eps))/(2\eta\mu)$ steps to reach it if $\eta<\eta^*$.
\end{enumerate}
Hence every constant step needs at least
\begin{equation}\label{eq:warmup:lowerconst}
  T\ge\Bigl(\frac{(K_0+K_1\Delta_0/2)\tanh a}{4\mu a}-\frac12\Bigr)\log\frac{\Delta_0}{2\eps}
\end{equation}
steps, where $a=\operatorname{arccosh}(1+K_1\Delta_0/(2K_0))\approx\log(2+K_1\Delta_0/K_0)$.
\end{proposition}

\begin{proof}
The Hessian is diagonal with entries $\mu\le K_0$ and $K_0\cosh(\sqrt{K_1}z)=K_0+K_1\Loss_{\mathrm c}(z)$, both at most $K_0+K_1\Loss$. For the gradient, $\Loss_{\mathrm c}'(z)^2=(K_0^2/K_1)\sinh^2(\sqrt{K_1}z)$ and $\sinh^2=(\cosh-1)(\cosh+1)\ge2(\cosh-1)$ give $\Loss_{\mathrm c}'^2\ge2K_0\Loss_{\mathrm c}\ge2\mu\Loss_{\mathrm c}$, and $(\mu x)^2=2\mu\cdot\frac\mu2x^2$; adding, $\norm{\nabla\Loss}_2^2\ge2\mu\Loss$.

Gradient descent acts on the two coordinates separately: $x_{t+1}=(1-\eta\mu)x_t$ and $z_{t+1}=\psi(z_t)$ with the odd map $\psi(z)=z-\eta(K_0/\sqrt{K_1})\sinh(\sqrt{K_1}z)$. For $z>0$, $\abs{\psi(z)}\ge z$ holds exactly when $\psi(z)\le-z$, that is, when $\eta\ge\eta^*(z)=2\sqrt{K_1}z/(K_0\sinh(\sqrt{K_1}z))$, and $\eta^*(z)$ decreases in $z$ because $b/\sinh b$ decreases on $(0,\infty)$. For (i), suppose that $\eta\ge\eta^*=\eta^*(z_0)$ and $\abs{z_t}\ge z_0$; then $\eta\ge\eta^*(\abs{z_t})$, so $\abs{z_{t+1}}=\abs{\psi(\abs{z_t})}\ge\abs{z_t}\ge z_0$. By induction $\Delta_t\ge\Loss_{\mathrm c}(z_t)\ge\Loss_{\mathrm c}(z_0)=\Delta_0/2>\eps$ for every $t$. For (ii), $\Delta_t\ge\frac\mu2x_t^2=(1-\eta\mu)^{2t}\Delta_0/2$. If $\eta\mu\ge1$ the claimed bound is not positive. Otherwise $\Delta_T\le\eps$ forces $2T\log(1/(1-\eta\mu))\ge\log(\Delta_0/(2\eps))$, and $\log(1/(1-s))\le s/(1-s)$ gives the claim. Finally, $(1-\eta\mu)/(2\eta\mu)=1/(2\eta\mu)-1/2$ decreases in $\eta$, so for $\eta<\eta^*$ it exceeds $1/(2\eta^*\mu)-1/2$, and $K_0\sinh a=K_0\cosh a\tanh a=(K_0+K_1\Delta_0/2)\tanh a$ gives \eqref{eq:warmup:lowerconst}.
\end{proof}

For $K_1\Delta_0\gg K_0$ the bound \eqref{eq:warmup:lowerconst} is about $K_1\Delta_0\log(\Delta_0/(2\eps))/(8\mu a)$, against the guarantee $(K_0\log(\Delta_0/\eps)+K_1\Delta_0)/\mu$ of the warmup: up to the logarithm $a$ and constants, a constant step pays the initial excess curvature once per factor $e$ of accuracy, and the warmup pays it once. The logarithm is real. The cheapest constant step that diverges on $\Loss_{\mathrm c}$ is $\eta^*=(2a\coth a)/(K_0+K_1\Delta_0/2)$, larger than the safe step by the factor $a\coth a$, because the curvature of $\Loss_{\mathrm c}$ falls quickly along a long step, and a step overshoots only if it exceeds the average curvature, not the curvature at its start. Constants matter as well: with $K_0=K_1=1$, $\mu=0.02$, $\Delta_0=50$ and $\eps=10^{-6}$, the best constant step on this instance needs $1398$ steps, \eqref{eq:warmup:lowerconst} gives $1392$, the loss-indexed warmup takes $1051$ steps, and \cref{thm:warmup:schedule} guarantees only $3386$. The comparison of worst cases exhibits the structure; the warmup model exhibits the size.

\Textcite{AlimisisIslamovLucchi2026}, who write $H_0,H_1$ for $K_0,K_1$, prove bounds of both kinds. For gradient descent with a warmup step size, under an aiming condition (a weakening of convexity) in place of \cref{ass:warmup:pl}, the term of their upper bound that grows as $\eps\to0$ involves $H_0$ alone, and $H_1$ enters only an additive term. For constant steps they prove lower bounds, among them, for functions satisfying \cref{ass:warmup:pl}, a number of iterations of order $(H_1/(4\mu))\,\Delta_0\log(\Delta_0/\eps)/(1+\log\Delta_0)$ in their normalization, with a hard function that is, like ours, quadratic near the minimum with an exponential tail; the logarithm of $\Delta_0$ in the denominator is the factor $a$ above. The next proposition makes the comparison sharp on the warmup model, where the stiff curvature is driven by the slow coordinate's loss.

\begin{proposition}[Warmup is optimal among stable schedules]\label[proposition]{prop:warmup:stable}
Consider the warmup model on the line $z=0$, with $\Delta_0=q(x_0)>0$, $S_t=K_0+K_1\Delta_t$ and $\Phi(D)=K_0\log D+K_1D$, and let $0<\eps<\Delta_0$.
\begin{enumerate}
\item Every schedule with $\eta_tS_t\le2$ for all $t$, so that no perturbation of the stiff coordinate is ever amplified, needs at least $(1-2\mu/K_0)^2\bigl(\Phi(\Delta_0)-\Phi(\eps)\bigr)/(4\mu)$ steps to reach $\Delta_T\le\eps$.
\item The loss-indexed schedule $\eta_t=c/S_t$ with $0<c\le2$ reaches it within $\bigl(\Phi(\Delta_0)-\Phi(\eps)\bigr)/\bigl(c\mu(2-c\mu/K_0)\bigr)$ steps.
\item Every constant step with $\eta S_0\le2$ needs at least $(1-2\mu/S_0)\,S_0\log(\Delta_0/\eps)/(4\mu)$ steps.
\end{enumerate}
\end{proposition}

\begin{proof}
By \cref{prop:warmup:model}(iii), $\Delta_{t+1}=(1-\mu\eta_t)^2\Delta_t$. For (i), $\mu\eta_t\le2\mu/S_t\le2\mu/K_0\le1$, and the inequality $\mu\eta_t<1$ is strict when $\Delta_t>0$, because then $S_t>K_0$; so the gaps stay positive, and $\Delta_t-\Delta_{t+1}=\Delta_t\mu\eta_t(2-\mu\eta_t)\le4\mu\Delta_t/S_t$. Since $\Phi'(D)=(K_0+K_1D)/D$ is decreasing,
\[
  \Phi(\Delta_t)-\Phi(\Delta_{t+1})\le(\Delta_t-\Delta_{t+1})\frac{S_{t+1}}{\Delta_{t+1}}\le\frac{4\mu\Delta_t}{\Delta_{t+1}}=\frac{4\mu}{(1-\mu\eta_t)^2}\le\frac{4\mu}{(1-2\mu/K_0)^2},
\]
and summing over $t<T$ with $\Phi(\Delta_T)\le\Phi(\eps)$ gives (i). For (ii), $c\mu/S_t\le2\mu/K_0\le1$ and
\[
  \Delta_t-\Delta_{t+1}=\Delta_t\frac{c\mu}{S_t}\Bigl(2-\frac{c\mu}{S_t}\Bigr)\ge c\mu\Bigl(2-\frac{c\mu}{K_0}\Bigr)\frac{\Delta_t}{S_t},
\]
so the integral comparison of part (iii) of \cref{thm:warmup:schedule} gives $\Phi(\Delta_t)\le\Phi(\Delta_0)-c\mu(2-c\mu/K_0)\,t$. For (iii), $\Delta_T=(1-\mu\eta)^{2T}\Delta_0\ge(1-2\mu/S_0)^{2T}\Delta_0$, and the inequality $\log(1/(1-s))\le s/(1-s)$ with $s=2\mu/S_0$ finishes as in \cref{prop:warmup:constant}.
\end{proof}

On the warmup model the theorem's schedule ($c=1$) is within a factor $2+O(1/\kappa)$ of every stable schedule, and riding the stability limit ($c=2$) is optimal among them up to $1+O(1/\kappa)$. The best stable constant step is slower by about $S_0\log(\Delta_0/\eps)/(\Phi(\Delta_0)-\Phi(\eps))$, which is close to $\min\{1+u_0,\log(\Delta_0/\eps)\}$. The guarantee of \cref{thm:warmup:schedule} is conservative here by the factor $2$ of the descent lemma, which certifies only half of the first-order decrease. Whether unstable schedules can do better is a different question, and \cref{sec:warmup:reading} answers it: they can.

\section{The theory: cross-entropy, normalized steps, Adam and the edge}\label{sec:warmup:more}

\subsection{Where loss-dependent curvature comes from}\label{sec:warmup:ce}

Losses built from exponentials have curvature bounded by their value. The exponential $\Loss(\theta)=e^{-\theta}$ on $\R$ has $\Loss''=\Loss$, so $K_0=0$ and $K_1=1$, and its infimum $0$ is not attained; \cref{ass:warmup:pl} fails, since $\Loss'^2=\Loss^2$. Still, the loss-indexed step $1/(K_1\Delta_t)=e^{\theta_t}$ gives $\theta_{t+1}=\theta_t+1$: the loss falls by the factor $e$ at every step, and the step size grows without bound, a warmup that never ends (Exercise~\ref{ex:warmup:exponential}). The cosh function \eqref{eq:warmup:cosh} has $\Loss_{\mathrm c}'^2=2K_0\Loss_{\mathrm c}+K_1\Loss_{\mathrm c}^2$, so it satisfies \cref{ass:warmup:pl} with $\mu=K_0$. For cross-entropy the bound holds with a constant that the plan for this chapter had wrong by a factor of two.

\begin{proposition}[Softmax regression]\label[proposition]{prop:warmup:softmax}
For an input $a\in\R^{d_{\mathrm{in}}}$, a class $c\in\{1,\dots,K\}$ and weights $W\in\R^{K\times d_{\mathrm{in}}}$, let $p=\softmax(Wa)$ and $\ell(W)=-\log p_c$. Then
\[
  \lmax\bigl(\nabla^2_W\ell\bigr)=\norm{a}_2^2\,\lmax\bigl(\diag(p)-pp^\top\bigr)\le2\norm{a}_2^2\,(1-p_c)\le2\norm{a}_2^2\,\ell(W),
\]
and the constant $2$ cannot be lowered. Consequently the average $\Loss$ of $\ell$ over a data set with $\norm{a_i}_2\le R$ satisfies \cref{ass:warmup:curvature} with $K_0=2R^2\Loss^\star$ and $K_1=2R^2$, and with $K_0=0$ when the data are separable.
\end{proposition}

\begin{proof}
With $s=Wa$, $\ell=-s_c+\log\sum_ke^{s_k}$ has gradient $p-e_c$ and Hessian $M=\diag(p)-pp^\top$ in $s$. Since $s$ is linear in $W$, the second derivative of $\ell$ in a direction $V\in\R^{K\times d_{\mathrm{in}}}$ is $(Va)^\top M(Va)$. As $V$ ranges over matrices with $\norm{V}_F=1$, the vector $Va$ ranges over the ball of radius $\norm{a}_2$, the extremes being $V=va^\top/\norm{a}_2$ with $\norm{v}_2=1$; so the largest eigenvalue is $\norm{a}_2^2\lmax(M)$. By Gershgorin's theorem every eigenvalue of $M$ is at most $\max_i(M_{ii}+\sum_{j\ne i}\abs{M_{ij}})=\max_i2p_i(1-p_i)$. For $i=c$ this is at most $2(1-p_c)$; for $i\ne c$ it is at most $2p_i\le2(1-p_c)$. Finally $1-p_c\le-\log p_c$. For $K=2$, $M$ has the eigenvalue $2p_c(1-p_c)$, whose ratio to $\ell=-\log p_c$ tends to $2$ as $p_c\to1$. Averaging over the data gives $\lmax(\nabla^2\Loss)\le2R^2\Loss=2R^2\Loss^\star+2R^2\Delta$; for separable data $\Loss^\star=0$.
\end{proof}

The plan for this chapter stated the bound with $1-p_c$ in place of $2(1-p_c)$, which fails for two classes as soon as $p_c>1/2$. For binary logistic regression with a single weight vector, $\ell(w)=\log(1+e^{-y\ip{w,a}})$, the Hessian is $\sigma(m)\sigma(-m)aa^\top$ with $m=y\ip{w,a}$ and $\sigma$ the logistic function, and the constant is $1$ (Exercise~\ref{ex:warmup:softmax}); the factor $2$ comes from the two rows of $W$ moving in opposite directions. So cross-entropy on a linear model, or on the last layer of a network with its features frozen, has curvature bounded by the loss, and when the data are separable, the step that \eqref{eq:warmup:gd} certifies grows without bound as the loss falls. The rest of a network is a different matter.

\begin{remark}[Rescaling symmetries break the global inequality]\label[remark]{rem:warmup:symmetry}
For the product model $\Loss(u,v)=\frac12(uv-y_\star)^2$ with $y_\star>0$, every point of the hyperbola $uv=y_\star$ is a global minimizer, with sharpness $u^2+v^2$, which is unbounded on it. Every sublevel set with positive gap contains the hyperbola, so no constants satisfy \cref{ass:warmup:curvature}. The same holds, with the square loss, for every network whose loss is invariant under multiplying one layer by $\alpha$ and the next by $1/\alpha$, as linear and ReLU networks are: their minimizers come in families along which the sharpness is unbounded. In such networks the assumption can hold only on a bounded region, with constants that depend on it, or along a trajectory.
\end{remark}

The product model is the subject of volume 2 of this series, \emph{The Hydrogen Atoms of Deep Learning}, Chapter 15 \parencite{HydrogenAtoms2026}, where gradient descent with a large step leaves a sharp minimizer through a catapult: the dynamics selects the curvature, which no inequality between curvature and loss can express. \Textcite{AlimisisIslamovLucchi2026} report that the condition is satisfied by common architectures, and \textcite{RiabininEtAl2026} check it along training trajectories; both are statements about the region that training visits.

\subsection{Normalized steps}\label{sec:warmup:normalized}

The theorem is about the step size of gradient descent, which multiplies the gradient. Optimizers whose update has a fixed length in some norm, such as sign descent, Muon and, roughly, Adam (\cref{ch:steepest,ch:sign,ch:spectral}), have a learning rate that is a length, and the question is whether that length must warm up. Let $\norm{\cdot}$ be a norm with dual norm $\norm{\cdot}_*$, and let normalized steepest descent take the step $\theta_{t+1}=\theta_t+\eta_t\lmo(\nabla\Loss(\theta_t))$ of length $\eta_t$ (\cref{ch:steepest}).

\begin{proposition}[When normalized steps need warmup]\label[proposition]{prop:warmup:normalized}
Suppose that on $\Theta_0$ the loss satisfies $u^\top\nabla^2\Loss(\theta)u\le K(\Delta(\theta))\norm{u}^2$ for every $u$, with $K\colon[0,\infty)\to(0,\infty)$ nondecreasing, and $\norm{\nabla\Loss(\theta)}_*\ge c\,\Delta(\theta)^\nu$ for some $c>0$ and $\nu\in(0,1]$. Then normalized steepest descent with the loss-indexed length $\eta_t=c\Delta_t^\nu/K(\Delta_t)$ satisfies
\[
  \Delta_{t+1}\le\Delta_t-\frac{c^2\Delta_t^{2\nu}}{2K(\Delta_t)},
\]
and since the gaps decrease, the length rises over a stretch of training exactly where $\Delta\mapsto\Delta^\nu/K(\Delta)$ decreases. For $K(\Delta)=K_0+K_1\Delta+K_\rho\Delta^\rho$ with $K_0>0$, $K_1,K_\rho\ge0$ and $\rho>1$:
\begin{enumerate}
\item if $\nu=1$ and $K_\rho=0$, the length only decreases: no warmup;
\item if $\nu=1$ and $K_\rho>0$, the length rises while $\Delta_t>\Delta'=\bigl(K_0/((\rho-1)K_\rho)\bigr)^{1/\rho}$ and decreases afterwards;
\item if $\nu=\frac12$ and $K_\rho=0$, the length rises while $K_1\Delta_t>K_0$ and decreases afterwards.
\end{enumerate}
\end{proposition}

\begin{proof}
Let $g=\nabla\Loss(\theta_t)$ and $u=\lmo(g)$, so that $\norm{u}=1$ and $\ip{g,u}=-\norm{g}_*$. If $\Delta_t=0$ the length is $0$; otherwise $g\ne0$. Points of the ray with $\Loss\le\Loss(\theta_t)$ have gap at most $\Delta_t$, and $K$ is nondecreasing, so \cref{lem:warmup:descent} applies with $M=K(\Delta_t)$. The length satisfies $\eta_t\le\norm{g}_*/K(\Delta_t)$, so
\[
  \Delta_{t+1}\le\Delta_t-\eta_t\norm{g}_*+\tfrac12K(\Delta_t)\eta_t^2\le\Delta_t-\tfrac12\eta_t\norm{g}_*\le\Delta_t-\tfrac12\eta_tc\Delta_t^\nu,
\]
which is the claim. The logarithmic derivative of $\Delta^\nu/K(\Delta)$ is $(\nu K-\Delta K')/(\Delta K)$, and $\nu K-\Delta K'$ equals $K_0$ in case (i), $K_0-(\rho-1)K_\rho\Delta^\rho$ in case (ii), and $\frac12(K_0-K_1\Delta)$ in case (iii).
\end{proof}

The gradient bound with $\nu=1$ holds for convex losses whose sublevel set $\Theta_0$ lies within distance $R$ of a minimizer, with $c=1/R$, because $\Delta\le\ip{\nabla\Loss(\theta),\theta-\theta^\star}\le R\norm{\nabla\Loss(\theta)}_*$; $\nu=\frac12$ is the Polyak--Łojasiewicz inequality in the dual norm. So under \cref{ass:warmup:curvature} a normalized method needs no warmup when the gradient's dual norm falls as fast as the loss, and it needs one when the curvature falls faster than the gradient. Case (ii) is the shape that \textcite{RiabininEtAl2026} derive for norm-constrained optimizers such as Muon and Lion: in their analysis the learning rate warms up and then decays only when $\rho>1$, with the turning point at the same gap $\Delta'$, which is where $\Delta/K(\Delta)$ is largest. In every case the length tends to $0$ with the gap: a normalized step must be decayed to converge, the thread of \cref{ch:momentum,ch:schedules}. The step length of gradient descent itself, $\eta_t\norm{\nabla\Loss(\theta_t)}_2$, follows case (iii) on the warmup model.

\subsection{Adam's first steps}\label{sec:warmup:adam}

Adam with bias correction \parencite{KingmaBa2015} keeps, from $m_{-1}=v_{-1}=0$, the moving averages $m_t=\beta_1m_{t-1}+(1-\beta_1)g_t$ and $v_t=\beta_2v_{t-1}+(1-\beta_2)g_t^2$ (coordinatewise), divides them by $1-\beta_1^{t+1}$ and $1-\beta_2^{t+1}$, and steps $\theta_{t+1}=\theta_t-\eta\,\hat m_t/(\sqrt{\hat v_t}+\epsilon)$. The \term{bias correction}\index{bias correction} removes the pull of the zero initial state, and it has a consequence for the first steps.

\begin{proposition}[Adam's first steps]\label[proposition]{prop:warmup:adamfirst}
Let $\epsilon=0$, with $0/0=0$, and fix a coordinate.
\begin{enumerate}
\item The first step is a sign step: $\theta_1-\theta_0=-\eta\sign(g_0)$, whatever the size of $g_0$.
\item For every $t\ge0$, $\hat v_t=\sum_{s=0}^tw_{t,s}g_s^2$ with weights $w_{t,s}=(1-\beta_2)\beta_2^{t-s}/(1-\beta_2^{t+1})$ that sum to $1$ and satisfy $\beta_2^t/(t+1)\le w_{t,s}\le\beta_2^{-t}/(t+1)$. If $t(1-\beta_2)\le\delta<1$, then $(1-\delta)\,\bar v_t\le\hat v_t\le\bar v_t/(1-\delta)$, where $\bar v_t=\frac1{t+1}\sum_{s\le t}g_s^2$ is the plain average of the squared gradients so far.
\item The effective number of averaged gradients is
\begin{equation}\label{eq:warmup:neff}
  n_t=\Bigl(\sum_{s=0}^tw_{t,s}^2\Bigr)^{-1}=\frac{1+\beta_2}{1-\beta_2}\cdot\frac{1-\beta_2^{t+1}}{1+\beta_2^{t+1}},
\end{equation}
which increases from $n_0=1$ to $(1+\beta_2)/(1-\beta_2)=2/(1-\beta_2)-1$. If the $g_s^2$ are independent with common mean $\sigma^2$ and variance $\varsigma^2$, then $\E\hat v_t=\sigma^2$ and $\Var(\hat v_t)=\varsigma^2/n_t$.
\end{enumerate}
\end{proposition}

\begin{proof}
At $t=0$, $\hat m_0=g_0$ and $\hat v_0=g_0^2$, so the step is $-\eta g_0/\abs{g_0}=-\eta\sign(g_0)$, which proves (i). Unrolling the recursion, $v_t=(1-\beta_2)\sum_{s=0}^t\beta_2^{t-s}g_s^2$ and $\sum_{k=0}^t\beta_2^k=(1-\beta_2^{t+1})/(1-\beta_2)$, which gives the weights and their sum. Each numerator $\beta_2^{t-s}$ lies in $[\beta_2^t,1]$ and the denominator $\sum_{k\le t}\beta_2^k$ in $[(t+1)\beta_2^t,t+1]$, which gives the bounds; Bernoulli's inequality $\beta_2^t\ge1-t(1-\beta_2)\ge1-\delta$ gives the comparison with $\bar v_t$. For (iii), $\sum_sw_{t,s}^2=(1-\beta_2)^2(1-\beta_2^{2t+2})/\bigl((1-\beta_2^2)(1-\beta_2^{t+1})^2\bigr)$, which simplifies to the reciprocal of \eqref{eq:warmup:neff}; the mean and variance follow from linearity and independence.
\end{proof}

So the first update has size $\eta$ in every coordinate, however small the gradient (\cref{ch:sign}), and for the first $\delta/(1-\beta_2)$ steps the preconditioner is the root mean square of the handful of gradients seen so far. With $\beta_1=\beta_2$ the bound of \cref{ch:sign} keeps every step at most $\eta$ per coordinate, but Adam's step is never small just because the gradient is: an optimizer that starts at its peak learning rate moves every coordinate by $\eta$ at once. RAdam replaces the first steps by a rectified version.

\begin{fieldguide}{RAdam}{\textcite{LiuEtAl2020RAdam}}
\fgrow{update}{Adam's $m,v$; after $k$ gradients, with $\rho_\infty=2/(1-\beta_2)-1$ and $\rho_k=\rho_\infty-2k\beta_2^k/(1-\beta_2^k)$, step $-\eta\,r_k\hat m/\sqrt{\hat v}$ if $\rho_k>4$ and $-\eta\hat m$ otherwise, $r_k=\bigl((\rho_k-4)(\rho_k-2)\rho_\infty/((\rho_\infty-4)(\rho_\infty-2)\rho_k)\bigr)^{1/2}$}
\fgrow{norm}{Euclidean for the first four steps (momentum SGD), then as Adam}
\fgrow{preconditioner}{$\diag(\sqrt{\hat v})$, damped by $r_k\approx(k(1-\beta_2)/2)^{1/2}$ while $5\ll k\ll1/(1-\beta_2)$}
\fgrow{averages}{as Adam; $\rho_k$ is the length of a simple moving average that matches the bias-corrected one, and $\rho_\infty=n_\infty$ of \eqref{eq:warmup:neff}}
\fgrow{state}{two numbers per parameter}
\fgrow{in the wild}{proposed to make warmup unnecessary; an untuned linear warmup over $2/(1-\beta_2)$ steps does as well \parencite{MaYarats2021}}
\end{fieldguide}

\begin{heuristic}[Heuristic derivation: the $2/(1-\beta_2)$ rule]
Suppose that the gradients of a coordinate in the first steps are independent draws from one distribution, and that a step is safe once the preconditioner is nearly as reliable as it will ever be. By \cref{prop:warmup:adamfirst}(iii) the relative variance of $\hat v_t$ is proportional to $1/n_t$, and
\[
  \frac{n_t}{n_\infty}=\tanh\Bigl(\frac{(t+1)\log(1/\beta_2)}2\Bigr),
\]
which reaches $\tanh1\approx0.76$ when $t+1=2/\log(1/\beta_2)\approx2/(1-\beta_2)$. A warmup that holds the step down until then lasts $2/(1-\beta_2)$ steps, the rule of thumb of \textcite{MaYarats2021}; RAdam's factor $r_k\approx\sqrt{k(1-\beta_2)/2}$ is a square-root warmup on the same time scale.
\end{heuristic}

The $2/(1-\beta_2)$ time scale is the \term{effective sample size}\index{effective sample size} of the second-moment estimate, the same number that sets the memory of Adam's preconditioner in \cref{ch:adaptive}. It is $2000$ steps for $\beta_2=0.999$ and $40$ steps for $\beta_2=0.95$. LLaMA uses $\beta_2=0.95$ and warms up for $2000$ steps: the long warmups of language-model recipes are not explained by the statistics of Adam's second moment alone.

\subsection{Warmup at the edge of stability}\label{sec:warmup:edge}

\Cref{ass:warmup:curvature} makes the curvature a function of the loss. In networks trained by gradient descent the causality also runs the other way: the sharpness rises to $2/\eta$ and stays there (\cref{ch:eos}). What does a rising $\eta_t$ do to it?

\begin{heuristic}[Heuristic derivation: the sharpness follows $2/\eta_t$]
Suppose that full-batch gradient descent with a slowly increasing schedule has reached the \term{edge of stability}\index{edge of stability}, $S(\theta_t)=2/\eta_t$, and that from then on the iterates follow the \term{central flow}\index{central flow} of \cref{ch:eos}: in the time $\tau=\sum_s\eta_s$,
\[
  \frac{d\theta}{d\tau}=-\nabla\Loss-\tfrac12\sigma^2\nabla S,
\]
with $\sigma^2$ the mean square of the oscillation along the top eigenvector of the Hessian, chosen so that the flow stays on the edge. With $\dot\eta=\eta_{t+1}-\eta_t$, the edge moves at the rate $d(2/\eta)/d\tau=-2\dot\eta/\eta^3$, and the flow moves the sharpness at the rate $\ip{\nabla S,d\theta/d\tau}=\alpha-\frac12\sigma^2\norm{\nabla S}_2^2$, with $\alpha=-\ip{\nabla\Loss,\nabla S}$ the rate of progressive sharpening. Equating the two,
\[
  \sigma^2=\frac{2\,(\alpha+2\dot\eta/\eta^3)}{\norm{\nabla S}_2^2}.
\]
The edge holds as long as $\sigma^2>0$, that is, as long as the threshold $2/\eta_t$ falls faster than the sharpness would fall under gradient flow, and then the sharpness follows $2/\eta_t$ down: the schedule, not the loss, sets the curvature. The oscillation carries an excess loss of about $\frac12S\sigma^2=\sigma^2/\eta$, which grows with the speed of the warmup. If $\eta_t$ rises so fast that $\sigma^2$ is not small, the cubic expansion behind the central flow fails, and the iterates catapult or diverge. On the warmup model $\alpha=-2\mu(S-K_0)$ and $\norm{\nabla S}_2^2=2\mu K_1(S-K_0)$, so on the edge $S=2/\eta$ the condition $\sigma^2>0$ reads $\dot\eta>\mu\eta^2(2-\eta K_0)$.
\end{heuristic}

If $\alpha>0$, the edge, once reached, holds for every nondecreasing schedule, and after the warmup the sharpness stays at $2/\eta_{\max}$ as in \cref{ch:eos}. If $\alpha<0$, as on the warmup model, it holds only while the warmup outruns the flattening that gradient descent produces on its own. In both cases the warmup lowers the curvature that the run will have at its peak learning rate, which is \cref{obs:warmup:sharpness}. The idealizations agree on the direction and disagree on the mechanism: in \cref{thm:warmup:schedule} the step waits for the loss to lower the curvature; at the edge the step lowers the curvature itself, through the oscillation, and the loss pays for it.

\section{Reading the theory}\label{sec:warmup:reading}

\paragraph{The shape and the length of the warmup.} \Cref{cor:warmup:profile} predicts a step size that starts at $1/(K_0+K_1\Delta_0)$, rises slowly along a hyperbola while the excess curvature dominates, turns concave at $2/(3K_0)$ and saturates exponentially at $1/K_0$; its length is $T_{1/2}=(K_0/\mu)(u_0-1+\log u_0)$. On the warmup model the rate $\mu$ is replaced by $2\mu$, because there the descent lemma is loose by its factor $2$. \Cref{fig:warmup:model} (right) compares this profile, for $K_0=K_1=1$, $\mu=0.02$ and $\Delta_0=50$, with gradient descent on the full model: the step reaches $1/(2K_0)$ at step $1324$, where the continuum predicts $1323$.

\begin{figure}[t]
\begin{fullwidth}
\centering
\begin{tikzpicture}
\begin{groupplot}[group style={group size=2 by 1,horizontal sep=17mm},wplot,width=82mm,height=52mm,xmin=0]
\nextgroupplot[xlabel={step $t$},ylabel={gap $\Delta_t$},ymode=log,ymin=1e-8,ymax=3e4,xmax=4000,ytick={1e-8,1e-4,1,1e4},legend style={at={(0.5,1.03)},anchor=south,font=\scriptsize},legend columns=2]
\addplot[theory] table[x=t,y=adrec] {figures/data/09-warmup-model.dat};
\addplot[simulation,each nth point=4,filter discard warning=false] table[x=t,y=adsim] {figures/data/09-warmup-model.dat};
\addplot[theory3] table[x=t,y=corec] {figures/data/09-warmup-model.dat};
\addplot[simulation2,each nth point=4,filter discard warning=false] table[x=t,y=cosim] {figures/data/09-warmup-model.dat};
\addplot[theory2] table[x=t,y=lcrec] {figures/data/09-warmup-model.dat};
\addplot[only marks,mark=triangle*,mark size=1.6pt,color=sea,mark options={fill=sea,draw=white,line width=0.2pt},each nth point=2,filter discard warning=false] table[x=t,y=lcsim] {figures/data/09-warmup-model.dat};
\addplot[guide] table[x=t,y=guarantee] {figures/data/09-warmup-model.dat};
\legend{loss-indexed: recursion,{\quad gradient descent},{constant $1/S_0$: recursion},{\quad gradient descent},{linear, $T_{\mathrm w}=2000$: recursion},{\quad gradient descent},{guarantee of \cref{thm:warmup:schedule}}}
\nextgroupplot[xlabel={step $t$},ylabel={step size $\eta_t$},ymode=log,ymin=5e-3,ymax=4,xmax=7000,legend style={at={(0.5,1.03)},anchor=south,font=\scriptsize},legend columns=2]
\addplot[theory] table[x=t,y=etacont] {figures/data/09-warmup-model.dat};
\addplot[simulation,each nth point=6,filter discard warning=false] table[x=t,y=etaad] {figures/data/09-warmup-model.dat};
\addplot[guide] table[x=t,y=limad] {figures/data/09-warmup-model.dat};
\addplot[theory3] table[x=t,y=etals] {figures/data/09-warmup-model.dat};
\addplot[theory2] table[x=t,y=limls] {figures/data/09-warmup-model.dat};
\legend{{profile, \cref{cor:warmup:profile} with $2\mu$},{loss-indexed, gradient descent},{its limit $2/S_t$},{linear, $T_{\mathrm w}=6150$},{its limit $2/S_t$}}
\end{groupplot}
\end{tikzpicture}
\end{fullwidth}
\caption{The warmup model with $K_0=K_1=1$, $\mu=0.02$ and $\Delta_0=50$ ($\kappa=u_0=50$), one slow and one stiff coordinate, the stiff one started at $z_0=10^{-6}$; no randomness (script \texttt{09-warmup-model.py}). Left: the gap under the loss-indexed warmup $\eta_t=1/(K_0+K_1\Loss(\theta_t))$, under the constant step $1/S_0=1/(K_0+K_1\Delta_0)$, and under a linear warmup to $1/K_0$ over $2000$ steps. Curves: the recursion $\Delta_{t+1}=(1-\mu\eta_t)^2\Delta_t$ of \cref{prop:warmup:model}; points: gradient descent on the full model. The linear warmup crosses the edge $\eta_tS_t=2$, its stiff coordinate grows, and a catapult lifts the loss to $6\cdot10^3$ before it converges, faster than the recursion. Dashed: the guarantee \eqref{eq:warmup:time}, conservative by the factor $2$ of the descent lemma. Without warmup, $\eta=1/K_0$, the loss exceeds $10^{13}$ after six steps. Right: the step size of the loss-indexed warmup against the profile of \cref{cor:warmup:profile} (with $2\mu$ for $\mu$), and the shortest linear warmup to $1/K_0$ that never crosses the edge, which touches its stability limit near $t=400$.}
\label{fig:warmup:model}
\end{figure}

\paragraph{What the warmup buys.} On the same model the loss-indexed warmup reaches $\Delta\le10^{-8}$ in $1805$ steps, against $1808$ predicted by \cref{prop:warmup:stable}(ii) with the exact factor $(1-\mu\eta_t)^2$ and $3617$ guaranteed by \cref{thm:warmup:schedule}. The constant step $1/S_0$ needs $28469$ steps, and the best stable constant step, just below $2/S_0$, needs $14304$, against the lower bound $14224$ of \cref{prop:warmup:stable}(iii). The ratio, about $8$, is the factor $\min\{1+u_0,\log(\Delta_0/\eps)\}$ of the theory with its constants: $(1+u_0)\log(\Delta_0/\eps)/(2(\log(\Delta_0/\eps)+u_0))=7.9$. It grows with the excess curvature at initialization and with the accuracy required, and it is $1$ when the network does not start sharp.

\paragraph{A linear warmup is the wrong shape.} Practice uses straight lines from zero, and the model says how long a straight line must be to stay stable.

\begin{heuristic}[Heuristic derivation: the shortest stable linear warmup]
On the line $z=0$ of the warmup model take $\eta_t=t/(K_0T_{\mathrm w})$ for $t\le T_{\mathrm w}$, and replace the recursion by its continuum $\dot\Delta=-2\mu\eta_t\Delta$, so that $\Delta_t=\Delta_0\exp(-\mu t^2/(K_0T_{\mathrm w}))$. With $s=t/T_{\mathrm w}$ and $\chi=\mu T_{\mathrm w}/K_0$, the condition $\eta_tS_t\le2$ reads $s+u_0se^{-\chi s^2}\le2$ for $s\in[0,1]$. The second term peaks at $s=1/\sqrt{2\chi}$ with the value $u_0/\sqrt{2e\chi}$, and evaluating the condition there gives $(1+u_0/\sqrt e)/\sqrt{2\chi}\le2$, that is,
\[
  T_{\mathrm w}\ge\frac{K_0}{8\mu}\Bigl(1+\frac{u_0}{\sqrt e}\Bigr)^2.
\]
After the warmup the step is constant and $S_t$ decreases, so the condition holds there too as soon as $u_0e^{-\chi}\le1$, which is weaker.
\end{heuristic}

For $\kappa=u_0=50$ the formula gives $6133$ steps, and bisection on the exact recursion gives $6150$; for $u_0$ from $5$ to $200$ the two agree within $6\%$. A line from the origin must pass below the limit $2/S_t$, which stays near $2/S_0$ until the loss has fallen (\cref{fig:warmup:model}, right), so the stable length is quadratic in the initial excess $u_0$, while the profile of \cref{cor:warmup:profile} starts at the limit and is linear in $u_0$. The ratio of the two lengths is about $u_0/(4e)$, here $4.6$, and the stable linear warmup reaches $10^{-8}$ in $2617$ steps instead of $1805$.

\paragraph{Shorter warmups and catapults.} The theory says nothing about schedules that cross the edge, and the model shows that they can win. Linear warmups to $1/K_0$ over $1400$ steps or fewer diverge. Between $1600$ and $2600$ steps the run crosses the edge, a catapult lifts the loss to between $2\cdot10^3$ and $10^4$, and the run then converges faster than the stable recursion predicts: $865$ steps at $T_{\mathrm w}=1624$, half the steps of the loss-indexed warmup. From $3000$ steps on, the run follows the recursion, although it crosses the edge for every $T_{\mathrm w}<6150$: the factor $\abs{1-\eta_tS_t}$ is close to $0$ when $\eta_tS_t\approx1$, which crushes the stiff coordinate, and it has no time to grow back before the edge is left. The recursion describes the run only while the stiff coordinate is negligible, and that is where the agreement must fail.

\paragraph{The edge in a network.} \Cref{fig:warmup:mlp} trains a small network from a sharp initialization, $S_0\approx132$, to the peak learning rate $0.1$, whose threshold is $2/\eta=20$. Before the edge is reached, the sharpness falls slowly on its own. Once $2/\eta_t$ meets it, after about $200$ steps of the long warmup and $100$ of the short one, the sharpness follows $2/\eta_t$ down, a few percent above it, as \cref{sec:warmup:edge} predicts; after the warmup it stays between $19.4$ and $23$ for the long warmup and settles near $20$ for the short one. The short warmup meets the edge with a catapult: its loss rises from $1.3$ to $6$ before the sharpness drops. Warmups of $50$ or $100$ steps diverge, and $150$ steps suffice. Unlike the warmup model, the network holds the edge rather than overshooting it, because every step of a network feeds its top eigendirection, whereas the stiff coordinate of the model is fed only by its initial value.

\begin{figure}[t]
\begin{fullwidth}
\centering
\begin{tikzpicture}
\begin{groupplot}[group style={group size=2 by 1,horizontal sep=17mm},wplot,width=82mm,height=50mm,xmin=0,xmax=2000]
\nextgroupplot[xlabel={step $t$},ylabel={sharpness $S(\theta_t)$},ymode=log,ymin=10,ymax=400,legend style={at={(0.5,1.03)},anchor=south,font=\scriptsize},legend columns=2]
\addplot[theory] table[x=t,y=edgeB] {figures/data/09-warmup-mlp.dat};
\addplot[simulation2] table[x=t,y=SB] {figures/data/09-warmup-mlp.dat};
\addplot[theory3] table[x=t,y=edgeA] {figures/data/09-warmup-mlp.dat};
\addplot[simulation] table[x=t,y=SA] {figures/data/09-warmup-mlp.dat};
\legend{{$2/\eta_t$, $T_{\mathrm w}=1000$},{measured},{$2/\eta_t$, $T_{\mathrm w}=250$},{measured}}
\nextgroupplot[xlabel={step $t$},ylabel={training loss},ymode=log,ymin=5e-4,ymax=30,legend style={at={(0.5,1.03)},anchor=south,font=\scriptsize},legend columns=3]
\addplot[simulation2] table[x=t,y=LB] {figures/data/09-warmup-mlp.dat};
\addplot[simulation] table[x=t,y=LA] {figures/data/09-warmup-mlp.dat};
\addplot[only marks,mark=triangle*,mark size=1.7pt,color=sea,mark options={fill=sea,draw=white,line width=0.2pt}] table[x=t,y=loss] {figures/data/09-warmup-mlp-nowarm.dat};
\legend{{$T_{\mathrm w}=1000$},{$T_{\mathrm w}=250$},{no warmup}}
\end{groupplot}
\end{tikzpicture}
\end{fullwidth}
\caption{Warmup at the edge in a small network: a tanh network with three hidden layers of width $64$ and output multiplier $5$ (a sharp initialization), trained by full-batch gradient descent on the square loss with $256$ Gaussian inputs in $\R^{10}$ labelled by a random tanh teacher, with a linear warmup to $\eta=0.1$ over $250$ or $1000$ steps (seeds fixed in \texttt{09-warmup-mlp.py}). Left: the sharpness, by power iteration with Hessian--vector products every $20$ steps (points), and the prediction $2/\eta_t$ of \cref{sec:warmup:edge} (curves), which holds once the edge is reached. Right: the training loss; without warmup it exceeds $10^9$ after six steps.}
\label{fig:warmup:mlp}
\end{figure}

\section{What the theory misses}\label{sec:warmup:misses}

\paragraph{What was idealized away.} A global inequality between curvature and loss: in networks it is measured along trajectories, and \cref{rem:warmup:symmetry} shows that it cannot hold on a sublevel set when the network has a rescaling symmetry. The Polyak--Łojasiewicz inequality, which makes the rates computable and is at best local in a network. Full batches: gradient noise adds its own reason to start small, the one \textcite{GoyalEtAl2017} gave, and at small batches the edge sits lower \parencite{AndreyevBeneventano2024}. A fixed preconditioner: for Adam the relevant curvature is the preconditioned sharpness, whose threshold is near $38/\eta$ for $\beta_1=0.9$ \parencite{CohenEtAl2022}, and the preconditioner itself changes fastest in the first steps (\cref{prop:warmup:adamfirst}). Normalization layers and weight decay, which make the effective learning rate an angular speed that is far from its equilibrium early in training (\cref{ch:decay}). Each idealization removes a reason for warmup, so the length that the theory prescribes is a lower bound on what practice needs only in the curvature mechanism, and not in general.

\paragraph{Where it breaks.} Three findings contradict the picture of \cref{thm:warmup:schedule}. First, the causality runs the other way in networks: the sharpness at the end of a warmup is $2/\eta_{\max}$, set by the schedule and not by the loss \parencite{CohenEtAl2021,KalraBarkeshli2024}; in \cref{fig:warmup:mlp} both warmups end at $S\approx20$ with different losses, so a bound $S\le K_0+K_1\Delta$ fitted to one run would be wrong for the other. Second, instability is not failure. Constant steps above $2/S_0$ converge through a catapult in wide networks \parencite{LewkowyczEtAl2020} and in the product model (volume 2 of this series, Chapter 15); for logistic regression, large steps with a non-monotone loss converge faster than every stable step \parencite{WuEtAl2024}; and on the warmup model, a linear warmup that crosses the edge converges in half the steps of the stable optimum of \cref{prop:warmup:stable}. The lower bounds of \cref{prop:warmup:constant,prop:warmup:stable} concern schedules that never cross the edge; the theory has nothing to say about the others. Third, the mechanism need not be curvature at all. With Pre-LN or a suitable initialization, transformers train without warmup \parencite{XiongEtAl2020,HuangEtAl2020}; for Adam, initializing the second moment with the first squared gradient or controlling the angular update recovers much of warmup's benefit \parencite{KalraBarkeshli2024,KossonMessmerJaggi2024}, which points to Adam's early statistics and update sizes rather than to a curvature that only time can lower. Finally, practice fixes the length by rules of thumb: $2000$ steps in LLaMA, where $2/(1-\beta_2)=40$ and where $K_1\Delta_0/\mu$ has never been measured; the theory reproduces the shape of the practice, not its numbers.

\paragraph{Repairs and open problems.} \Textcite{AlimisisIslamovLucchi2026} turn their analysis into a warmup schedule with one parameter, which matches a tuned linear warmup on language and vision models; \textcite{RiabininEtAl2026} derive a schedule for norm-constrained optimizers that sets the warmup length automatically and, in LLaMA pretraining, matches or beats the best tuned warmups; \textcite{LiuEtAl2025Warmup}, with curvature bounded by $K_0+K_\rho\Delta^\rho$, show that warmup can make gradient descent up to $\Theta(T)$ times faster than every non-increasing schedule in some cases (Exercise~\ref{ex:warmup:power}). On the side of the optimizer, GI-Adam, Pre-LN, rescaled initializations and the update controls of \textcite{KossonMessmerJaggi2024} remove much of the need, each for its own reason. Open: a derivation of the curvature--loss inequality along the trajectories of a network, even a two-layer one (Exercise~\ref{ex:warmup:twolayer}); lower bounds over all schedules, including those that cross the edge (Exercise~\ref{ex:warmup:catapult}); a theory of warmup with gradient noise, where the edge itself depends on the batch size; and the preconditioned version of the whole story for Adam, in which the preconditioner and the curvature evolve on different time scales.

\section{Back in the wild}\label{sec:warmup:back}

\fieldnotes{Linear warmup of the learning rate from nearly zero over a few thousand steps ($2000$ for LLaMA, $375$ million tokens for GPT-3, five epochs for large-batch ImageNet).}
{Curvature bounded by the loss, $S\le K_0+K_1\Delta$, with the Polyak--Łojasiewicz inequality; full-batch gradient descent; the warmup model; Adam's moving averages; the central flow at the edge.}
{A local descent lemma; a loss-indexed or scheduled warmup reaches $\eps$ in $(K_0\log(\Delta_0/\eps)+K_1\Delta_0)/\mu$ steps, while a constant step pays the initial curvature at every scale; the exact profile; optimality among stable schedules on the warmup model; softmax curvature at most $2R^2\ell$; Adam's first step is a sign step.}
{A hyperbolic rise that saturates, of length about $K_1\Delta_0/\mu$; a stable linear warmup must be $u_0/(4e)$ times longer; the sharpness follows $2/\eta_t$ at the edge (derived); Adam's statistics settle after $2/(1-\beta_2)$ steps (derived).}
{The inequality is measured, not derived, and fails globally under rescaling symmetries; catapults let unstable warmups succeed and even win; noise, preconditioning and normalization layers.}
{Warm up until the curvature has fallen below what the peak step tolerates, and expect the needed length to grow with how sharp the initialization is (proved in the idealization). Indexing the warmup by the loss or by a measured sharpness, rather than by a fixed count, is suggested by the theory and by two recent schedulers; long warmups at $\beta_2=0.95$ are not explained by Adam's statistics (derived).}

The theory gives a practitioner one firm statement and one suggestion. The firm statement, proved in the idealization, is that a constant step must be safe at the start and then pays for the initial curvature at every scale of accuracy, while a warmup pays once; the gain is large only when the initialization is much sharper than the trained network, which a power iteration on the first minibatch can check against $2/\eta_{\max}$. The suggestion is the shape: start at the largest step that the current curvature allows and rise as it falls, which is a hyperbola in time, not a line from zero. The theory does not give the length of a language-model warmup, and it should not be read as forbidding short ones; at the edge, and through catapults, networks lower their own curvature faster than the stable theory allows.

The chapter sits between the edge of stability of \cref{ch:eos}, which supplies the mechanism of \cref{sec:warmup:edge}, and the clipping of \cref{ch:spikes}, which tames a curvature that grows with the gradient rather than with the loss. Warmup is the first phase of the schedules of \cref{ch:schedules}, and its interaction with normalized updates continues in the effective learning rate of \cref{ch:decay}: with weight decay and normalization layers, the angular update that \textcite{KossonMessmerJaggi2024} control early in training is the quantity that equilibrates late. That Adam's first step is a sign step connects it to \cref{ch:sign}, and the role of the norm in \cref{prop:warmup:normalized} to \cref{ch:steepest}.

\begin{exercises}
\item \consolidation (The softmax constant.) (a) Show that $\lmax(\diag(p)-pp^\top)\le2\max_ip_i(1-p_i)\le2(1-p_c)$ for every probability vector $p$ and class $c$. (b) For $K=2$ show that the top eigenvalue is $2p_c(1-p_c)$, and conclude that the constant $2$ of \cref{prop:warmup:softmax} cannot be lowered and that $\lmax\le1-p_c$ fails when $p_c>\frac12$. (c) For binary logistic regression with one weight vector, $\ell(w)=\log(1+e^{-y\ip{w,a}})$, show that $\nabla^2\ell=\sigma(m)\sigma(-m)\,aa^\top\preceq\ell(w)\,aa^\top$ with $m=y\ip{w,a}$.\label{ex:warmup:softmax}

\item \consolidation (The extremal function.) (a) Show that \cref{ass:warmup:curvature} implies $\norm{\nabla\Loss}_2^2\le2(K_0+K_1\Delta)\Delta$ on $\Theta_0$. (b) For $\Loss_{\mathrm c}$ of \eqref{eq:warmup:cosh} show that $\Loss_{\mathrm c}'^2=2K_0\Loss_{\mathrm c}+K_1\Loss_{\mathrm c}^2$. (c) Show that gradient descent with the loss-indexed step on $\Loss_{\mathrm c}$ satisfies $a_{t+1}=a_t-\tanh a_t$ for $a_t=\sqrt{K_1}z_t$, so that the loss falls by a factor close to $e$ per step while $a_t$ is large, and the warmup on this function lasts about $\log(K_1\Delta_0/K_0)$ steps rather than $K_1\Delta_0/\mu$. Which hypothesis of \cref{thm:warmup:schedule} is far from tight here?\label{ex:warmup:extremal}

\item \consolidation (The profile.) (a) Derive the formula of \cref{cor:warmup:profile} directly from $\dot D=-r(D)$. (b) Show that $\eta(t)$ is convex while $\eta(t)<2/(3K_0)$ and concave afterwards. \hint{with $u=K_1D/K_0$, $\dot u=-(\mu/K_0)u/(1+u)$, and $\ddot\eta$ has the sign of $2u-1$.} (c) For $K_0=K_1=1$, $\mu=0.01$ and $\Delta_0=100$, compute $T_{1/2}$ and the time at which $\eta$ reaches $0.9/K_0$.\label{ex:warmup:profile}

\item \consolidation\hard (The hypothesis $K_0\ge2\mu$.) Show that the warmup model satisfies $\lmax(\nabla^2\Loss)\le K_0+K_1\Loss$ everywhere if and only if $K_0\ge25\mu/16$. \hint{in the proof of \cref{prop:warmup:model} let $\varepsilon\to0$: the condition becomes $(S-\mu)S\ge4\mu K_1q$ for all $q\ge0$, a quadratic inequality in $K_1q$ whose discriminant is $25\mu^2-16K_0\mu$.}\label{ex:warmup:threshold}

\item \consolidation (Adam's effective sample size.) (a) Prove \eqref{eq:warmup:neff} and the formula $n_t/n_\infty=\tanh((t+1)\log(1/\beta_2)/2)$. (b) Show that RAdam's $\rho_k$ and $n_{k-1}$ both equal $k$ to first order in $k(1-\beta_2)$, and that both tend to $(1+\beta_2)/(1-\beta_2)$. (c) Show that with $\beta_1>0$ the first step is still $-\eta\sign(g_0)$, and compute the second step when $g_1=-g_0$.\label{ex:warmup:neffex}

\item \variation (No curvature floor: the exponential loss.) Let $\Loss(w)=\frac1n\sum_i\exp(-y_i\ip{w,a_i})$ with $\norm{a_i}_2\le R$, and suppose the data are separable with margin $\gamma$: some unit vector $\bar w$ has $y_i\ip{\bar w,a_i}\ge\gamma$ for all $i$. (a) Show that $\lmax(\nabla^2\Loss(w))\le R^2\Loss(w)$ and $\norm{\nabla\Loss(w)}_2\ge\gamma\Loss(w)$. (b) Show that gradient descent with the loss-indexed step $\eta_t=1/(R^2\Loss(w_t))$ satisfies $\Loss(w_{t+1})\le(1-\gamma^2/(2R^2))\Loss(w_t)$: linear convergence without the Polyak--Łojasiewicz inequality, with step sizes that grow geometrically. (c) Contrast with a constant step, under which the loss decays only polynomially (\cref{ch:bias}), and relate the result to the large-step regime of \textcite{WuEtAl2024}.\label{ex:warmup:exponential}

\item \variation (Curvature that grows like a power of the gap.) Suppose that $S\le K_0+K_\rho\Delta^\rho$ on $\Theta_0$ with $0<\rho\le1$, as for \textcite{LiuEtAl2025Warmup}, and that \cref{ass:warmup:pl} holds. Show that the loss-indexed step $1/(K_0+K_\rho\Delta_t^\rho)$ reaches $\eps$ within $\bigl(K_0\log(\Delta_0/\eps)+K_\rho(\Delta_0^\rho-\eps^\rho)/\rho\bigr)/\mu$ steps. Where does the proof of \cref{thm:warmup:schedule} use $\rho\le1$?\label{ex:warmup:power}

\item \variation (Momentum on the warmup model.) With the moving-average momentum of \cref{ch:momentum}, $m_t=\beta m_{t-1}+(1-\beta)\nabla\Loss(\theta_t)$ and $\theta_{t+1}=\theta_t-\eta m_t$, the stiff coordinate with frozen curvature $S$ is stable if and only if $\eta S<2(1+\beta)/(1-\beta)$. Redo the heuristic derivation of the shortest stable linear warmup with this threshold in place of $2$, assuming that the slow coordinate still obeys $\dot\Delta=-2\mu\eta_t\Delta$, and evaluate the result for $\beta=0.9$.\label{ex:warmup:momentum}

\item \variation (The step length of gradient descent.) On the line $z=0$ of the warmup model, show that the step length $\eta_t\norm{\nabla\Loss(\theta_t)}_2$ of the loss-indexed warmup is $\sqrt{2\mu\Delta_t}/(K_0+K_1\Delta_t)$, that it rises until $K_1\Delta_t=K_0$ and falls afterwards, and that this is case (iii) of \cref{prop:warmup:normalized}. What does a sign-descent run with a fixed step length do on this model when that length exceeds the peak?\label{ex:warmup:length}

\item \simulation (The warmup model.) With \texttt{09-warmup-model.py}: (a) check the shortest stable linear warmup against the heuristic formula for $u_0\in\{5,10,20,50,100,200\}$; (b) scan the warmup length and locate the window of catapults ($T_{\mathrm w}$ between about $1600$ and $2600$ for the default parameters), recording the largest loss of each run; (c) run the scheduled warmup of \cref{thm:warmup:schedule}(ii) and the loss-indexed one with $c=2$, and explain their step counts by \cref{prop:warmup:stable}.\label{ex:warmup:simmodel}

\item \simulation (A network at the edge.) With \texttt{09-warmup-mlp.py}: (a) measure $\eta_tS(\theta_t)$ during and after the warmup and check that it stays close to $2$ once the edge is reached; (b) find the shortest warmup that trains for the peak learning rates $0.05$ and $0.2$ (it is between $100$ and $150$ steps for $0.1$), and compare how it scales with the peak; (c) set the output multiplier to $1$, where $S_0\approx5.7$, and observe that with a peak whose threshold $2/\eta$ exceeds $S_0$ the sharpness rises toward $2/\eta_t$ instead of falling.\label{ex:warmup:simmlp}

\item \research (The inequality for a two-layer network.) \Cref{rem:warmup:symmetry} rules out \cref{ass:warmup:curvature} on sublevel sets of networks with rescaling symmetries. Find conditions, such as balanced or small initializations, bounded weight norms, or a restriction to gradient-descent trajectories, under which a two-layer network with the square or cross-entropy loss satisfies $S\le K_0+K_1\Delta$ with explicit constants, and compute $u_0=K_1\Delta_0/K_0$ for standard initializations. As far as we know, no such derivation exists.\label{ex:warmup:twolayer}

\item \research (Lower bounds that allow catapults.) On the warmup model, catapulting warmups beat every stable schedule (\cref{sec:warmup:reading}). Is there a lower bound over all schedules, stable or not, for gradient descent on the class of \cref{ass:warmup:curvature,ass:warmup:pl}? What is the fastest schedule on the warmup model, and does it cross the edge? Both questions are open.\label{ex:warmup:catapult}
\end{exercises}

\notesandrefs
Warmup entered deep learning as a fix for runs that would not start: \textcite{HeEtAl2016} held a $110$-layer residual network at a tenth of its learning rate for its first $400$ iterations, \textcite{GoyalEtAl2017} made the ramp linear to train with large minibatches, and \textcite{VaswaniEtAl2017} paired a linear warmup with an inverse-square-root decay. The recipes of GPT-3 and LLaMA \parencite{BrownEtAl2020,TouvronEtAl2023} fixed the modern form, a linear ramp followed by cosine decay. \Textcite{XiongEtAl2020} traced the need in Post-LN Transformers to large gradients near the output at initialization, and \textcite{HuangEtAl2020} removed it by initialization.

The curvature view of warmup is empirical: \textcite{GilmerEtAl2022} tied early instabilities to the sharpness, and \textcite{KalraBarkeshli2024} described the sharpness-reduction and progressive-sharpening regimes of warmup and proposed GI-Adam. The edge of stability is due to \textcite{CohenEtAl2021}, its self-stabilization to \textcite{DamianNichaniLee2023}, the central flow to \textcite{CohenEtAl2025}, the catapult to \textcite{LewkowyczEtAl2020}, and the solvable product model with its catapult to volume 2 of this series \parencite{HydrogenAtoms2026}. Large steps with a non-monotone loss on logistic regression are analysed by \textcite{WuEtAl2024}.

The theory of warmup under loss-dependent curvature is recent. \Textcite{AlimisisIslamovLucchi2026} introduced $(H_0,H_1)$-smoothness, proved upper bounds for gradient descent with a warmup step and lower bounds for constant steps, and turned the analysis into a one-parameter schedule. \Textcite{LiuEtAl2025Warmup} proved acceleration by warmup under a power-law version of the condition, and \textcite{RiabininEtAl2026}, in a preprint, extended the analysis to norm-constrained optimizers and derived an adaptive scheduler. The loss-dependent condition is a relative of the gradient-dependent $(L_0,L_1)$-smoothness of \textcite{ZhangHeSraJadbabaie2020}. The Polyak--Łojasiewicz inequality is surveyed by \textcite{KarimiNutiniSchmidt2016}. As far as we know, the warmup model, the profile of \cref{cor:warmup:profile} in terms of the Wright omega function, \cref{prop:warmup:stable} and the shortest stable linear warmup are new here.

For Adam, the bias correction is in \textcite{KingmaBa2015}; the variance argument and RAdam are due to \textcite{LiuEtAl2020RAdam}, the rule of thumb $2/(1-\beta_2)$ to \textcite{MaYarats2021}, and the angular-update analysis of warmup in GPT training to \textcite{KossonMessmerJaggi2024}. Adam's edge, with its threshold near $38/\eta$, is in \textcite{CohenEtAl2022}; the minibatch edge in \textcite{AndreyevBeneventano2024}. \Textcite{WortsmanEtAl2024} study warmup among the knobs of small proxy models for large-scale instabilities.

\paragraph{Translating the literature.} \Textcite{AlimisisIslamovLucchi2026} write $f$ for $\Loss$, $w$ for $\theta$, $f^*$ for $\Loss^\star$ and $H_0,H_1$ for $K_0,K_1$, and bound $\norm{\nabla^2f(w)}$ by $H_0+H_1(f(w)-f^*)$; their warmup step is a multiple $\theta\eta_k$ of a base step $\eta_k$, so their $\theta$ is a constant and not the parameter vector, and they count iterations by $K$. \Textcite{RiabininEtAl2026} write $x$ for $\theta$ and $\Delta^t$ for our $\Delta_t$, and bound the dual norm of gradient differences by $K(x)\norm{x-y}$ with $K(x)=K_0+K_1(f(x)-f^\star)+K_\rho(f(x)-f^\star)^\rho$; with $K_\rho=0$ and the Euclidean norm their assumption reduces to that of Alimisis et al. \Textcite{LiuEtAl2025Warmup} call their condition $(\rho,K_0,K_\rho)$-smoothness, a bound $K_0+K_\rho(f(w)-f^*)^\rho$ on the local curvature. In RAdam, $\rho_t$ is indexed by the number of gradients, our $t+1$.

\printchapterbib
\end{refsection}
